
sera obligé de travailler avec l'anneau de Chow, puisque c'est sur ce dernier que se trouve défini naturellement l'homomorphisme @@M0@@ comme homomorphisme à valeurs dans @@M1@@, comparer (XIV 4.1). Malheureusement, on n'a pas développé pour le moment une théorie de la spécialisation pour l'anneau de Chow (cf. 7.14.), qui permettrait d'écrire des carrés commutatifs analogues à (7.13.9) et (7.13.10) en termes des anneaux de Chow et de la cohomologie à coefficients @@M2@@. Tout ce qu'on peut dire lorsque @@M3@@ est lisse et propre, c'est qu'on a un carré commutatif
@@M4@@
analogue à (7.13.9), et un carré correspondant analogue à (7.13.10) qu'on laisse au lecteur le soin d'écrire, où @@M5@@ désigne le groupe des cycles algébriques (modulo rien du tout) de codimension @@M6@@, où les flèches horizontales sont les homomorphismes «~classe de cohomologie associée~» définis dans (SGA 5 IV), et où @@M7@@ désigne l'homomorphisme de spécialisation des cycles algébriques, associant à un cycle @@M8@@ le cycle @@M9@@, où @@M10@@ est l'adhérence schématique de @@M11@@ dans @@M12@@ (cf. 7.14). Cette commutativité résulte du fait suivant, établi dans (SGA 5 IV); si @@M13@@ est un sous-schéma fermé de @@M14@@ de codimension @@M15@@ sur chaque fibre, on peut lui associer une classe de cohomologie @@M16@@, dont la formation commute à tout changement de base, et en particulier telle que sa restriction à @@M17@@ (resp. @@M18@@) soit égale à @@M19@@ (resp. @@M20@@). Bien entendu, lorsqu'on sait que l'homomorphisme @@M21@@ dans (7.13.12) transforme cycles rationnellement équivalents à zéro en cycles rationnellement équivalents à zéro (@@M22@@ étant projectif sur @@M23@@),
